\chapter{Simplicity, Probability and Explanatory Preference}
\label{sec:p19-explanatory-preferences}

An explanation can be sufficient under a stipulated causal model,
preferred by a participant, or more probable under a specified law.
These are different claims. This chapter treats Tania Lombrozo's
complete four-experiment presentation, \emph{Simplicity and probability
in causal explanation} \cite{Lombrozo2007SimplicityProbability}.
The held article is the published \emph{Cognitive Psychology} 55
version, pages 232--257, dated 2007; it was accepted on 29 September
2006 and first available online on 9 November 2006.
Its operational measure of simplicity is the number of causes
invoked, not description length or number of fitted parameters.
The article is this chapter's bounded empirical corpus.

The task uses fictional aliens, diseases and symptoms whose causal
relations are provided to the participant. It tests evaluation of
supplied explanations for a particular case. It does not test discovery
of an unknown causal relation or spontaneous generation of every
candidate explanation. The evidence record follows
\cref{def:p19-evidence-record}; the mathematical comparisons below
are independently derived handbook constructions.

\section{The explanatory task and its logical choices}
\label{sec:p19-explanation-task}

\begin{hbdefinition}{def:p19-explanation-task}{One common cause and two separate causes}
Use \(D_1,D_2,D_3\) for the three disease propositions and
\(S_1,S_2\) for two symptoms. The stipulated descriptions say that
\(D_1\) always causes both symptoms, \(D_2\) always causes \(S_1\)
but does not cause \(S_2\), and \(D_3\) always causes \(S_2\)
but does not cause \(S_1\). Nothing else is known to cause the
symptoms. The target alien has both symptoms.

In the six-option task, the alternatives are each single disease
and each pair of diseases. The focal comparison is \(D_1\)
against \(D_2\) and \(D_3\). The former invokes one cause and
the latter two. \(D_2\) alone and \(D_3\) alone fail to account
for one symptom. \(D_1\) with either other disease still accounts
for both, but adds a cause to an already sufficient explanation.
There is no three-disease option in this menu.

The principal question asks for the \emph{most satisfying}
explanation. It elicits one categorical choice, not a numerical
probability or a satisfaction rating. A later question requests
a written justification. The terms satisfying and likely are
deliberately varied in Experiment 4.
\end{hbdefinition}

For example, the source's first scenario concerns a population of
750 aliens on Zorg, Treda's sore minttels and purple spots, and
Tritchet's syndrome, Morad's disease and a Humel infection.
One symptom is recognizable and the other invented, to reduce
participants' prior etiological expectations. The symptoms and
disease-role assignments are counterbalanced.
Saying that one disease does not \emph{cause} a symptom does
not say that it prevents another disease from producing it.

\begin{hbtheorem}{thm:p19-explanation-logical-menu}{proved=annex:p19.explanation-logical-menu}{A precise finite-abduction correspondence}
For the explicitly stronger, closed deterministic model
\[
 T=\{S_1\leftrightarrow(D_1\lor D_2),\
      S_2\leftrightarrow(D_1\lor D_3)\},
\]
take permitted hypotheses \(H=\{D_1,D_2,D_3\}\) and observations
\(M=\{S_1,S_2\}\) in \cref{def:p16-abduction-input}.
The subset-minimal explanations are exactly
\(\{D_1\}\) and \(\{D_2,D_3\}\); only \(\{D_1\}\) has minimum
cardinality. Both are consistent and sufficient.
Adding \(D_2\to\neg S_2\) would instead make the two-cause
candidate inconsistent. That extra rule is not the original
statement that \(D_2\) does not cause \(S_2\).
\end{hbtheorem}

This correspondence identifies two formal criteria on this declared
model. It does not establish that participants implement an abduction
algorithm, or that a minimum-cardinality explanation is most probable.
The distinction between logical sufficiency and truth of explanatory
premises also belongs to the D--N treatment in
\cref{def:p16-dn-schema}. Participants' preferences establish neither
the truth nor the law status required there.

\section{Experiment 1: separate simplicity and prevalence comparisons}
\label{sec:p19-explanation-experiment-one}

Forty-eight undergraduates and summer-school students from an
unnamed elite university participated for course credit, payment
or a small gift. The source reports 44 percent male, mean age 21,
and age standard deviation 3. Participants were randomly assigned
to two conditions; the method does not separately print each
condition's participant count.

The two-page questionnaire first presented the alien case and
explanation menu, then requested a justification.
The simplicity condition used the six-option task without disease
prevalences. The probability condition instead provided two
diseases, each always producing both symptoms, with approximately
50 and 73 cases respectively among 750 aliens. Its menu contained
either disease alone or both together. Thus the focal alternatives
had equal cause counts but different stated prevalences.
Four symptom sets were used. Disease order, the name assigned to
the common cause, and symptoms were counterbalanced, and answers
were presented in several random orders
\cite[pp.~236--237]{Lombrozo2007SimplicityProbability}.

In the simplicity condition 96 percent chose the one-cause
explanation. The article reports binomial \(p<.01\) both against
uniform choice among six options and against equal choice between
the two focal explanations. In the probability condition 92 percent
chose the more prevalent disease, again with \(p<.01\) against
uniform choice among three options and against the two-option
comparison. These are specified equal-choice nulls; they do not
give the probability that a participant actually used guessing.
Exact response counts are not supplied here by converting rounded
percentages into presumed condition sizes.

Two coders classified all 48 justifications, with 96 percent
initial agreement and disagreements resolved by discussion.
\emph{Simple} included explicit simplicity and related descriptions
such as reduced complexity; \emph{sufficient} invoked the one
disease's adequacy without explicitly naming simplicity;
\emph{probable} invoked greater likelihood; and \emph{other}
covered the remainder. Four were classified other.
Among those choosing the simpler explanation, the reported
percentages were 17 simple, 39 sufficient and 39 probable.
Every participant choosing the more prevalent explanation in
the probability condition appealed to probability.
Overall 44 of 48 justifications invoked probability, simplicity
or sufficiency \cite[pp.~237--238]{Lombrozo2007SimplicityProbability}.

The observed simplicity preference does not isolate its psychological
basis. It is compatible with valuing simplicity itself and with
using simplicity as a cue to probability. Written justifications
constrain those accounts but do not uniquely identify the process
producing a choice.

\section{Probability benchmarks need event definitions}
\label{sec:p19-explanation-probability}

\begin{hbdefinition}{def:p19-explanation-probability}{Inclusive explanations and an observation event}
On a finite probability space, let \(D_1,D_2,D_3\) now denote
events. Put \(H_s=D_1\), \(H_c=D_2\cap D_3\), and
\(E=H_s\cup H_c\), corresponding to both symptoms in the closed
deterministic model. Write
\[
 p_1=\mathbb P(H_s),\quad j=\mathbb P(H_c),\quad
 t=\mathbb P(H_s\cap H_c).
\]
Assume \(p_1+j-t>0\) when conditioning on \(E\).
The target is interpreted under this declared population law,
conditioned only on the two symptoms; an additional case-selection
mechanism would require its own conditioning model.
The inclusive meanings allow a case to satisfy both explanations:
asserting \(D_1\) makes no assertion that \(D_2,D_3\) are absent.
This is the interpretation expressly selected by the source's
note 4 \cite[p.~240]{Lombrozo2007SimplicityProbability}.
\end{hbdefinition}

\begin{hbtheorem}{thm:p19-explanation-probability}{proved=annex:p19.explanation-probability}{Overlap, marginal information and a worked comparison}
Under the displayed model,
\[
 \mathbb P(H_s\mid E)=\frac{p_1}{p_1+j-t},\qquad
 \mathbb P(H_c\mid E)=\frac{j}{p_1+j-t}.
\]
For \(j>0\), their ratio is \(R=p_1/j\).
The normalized quantity \(q_0=R/(1+R)=p_1/(p_1+j)\)
is not generally \(\mathbb P(H_s\mid E)\).
If \(p_1>0\), equality holds exactly when \(t=0\).
Under independence of \(D_2,D_3\), one may replace \(j\) by
\(\mathbb P(D_2)\mathbb P(D_3)\); the marginals alone do not
justify that replacement.

For independent disease indicators with prevalences
\((.1,.4,.4)\), the two posteriors are \(25/61\) and \(40/61\),
while \(q_0=5/13\).
Keeping the same three prevalences, but making \(D_2,D_3\)
mutually exclusive gives posteriors 1 and 0; making them
identical gives \(5/23\) and \(20/23\), with \(D_1\) independent
of that pair in each construction.
Thus the marginal prevalences alone do not determine the comparison.
\end{hbtheorem}

The posterior-odds identity of \cref{thm:p15-model-posterior-odds}
requires a declared model family, data laws and priors. A prevalence
ratio, an event posterior and a population choice proportion are
not interchangeable inputs to that identity.
The paper's Figure 1 labels \(R/(1+R)\) as a Bayesian-posterior
reference curve. With its inclusive event interpretation, this
curve can instead be read as normalized support for two answer
labels, with a further probability-matching response assumption.
Bayes' theorem alone does not require a person to choose the
more probable answer only with that probability: deterministic
maximization is another possible decision rule.

The source's note 2 invokes a disease-prevalence threshold of .5.
For inclusive events with equal prevalence \(p\) and independent
\(D_2,D_3\), however, the comparison is \(p\) against \(p^2\);
the former is larger for every \(0<p<1\).
The .5 threshold instead arises when comparing \emph{exclusive}
configurations \(D_1\) alone and \(D_2,D_3\) without \(D_1\)
under mutual independence. Their probabilities are
\(p(1-p)^2\) and \(p^2(1-p)\).
The annex proves both statements, preserving the event distinction.

\section{Experiment 2: eight prevalence conditions}
\label{sec:p19-explanation-experiment-two}

The retained sample was 144 undergraduate and summer-school
students from an unnamed elite university: 60 percent female,
mean age 22, standard deviation 5. Compensation was course credit,
payment or a small gift. Three additional participants were
replaced for blank responses and four for non-target explanation
choices. The latter justifications indicated a different intended
choice, such as naming a disease while claiming it caused both
symptoms. These replacements mean that the analyzed non-simple
choices are precisely the focal two-cause choices.

The questionnaire retained the six-option task and added each
disease's prevalence. There were eighteen participants randomly
assigned to each of eight conditions. The complete Table 1
frequencies below refer to the fictional population of 750:
\[
\begin{array}{c|rrr}
 \text{Nominal ratio}&D_1&D_2&D_3\\\hline
 15:1&50&50&50\\
 1:1&50&197&190\\
 9:10&50&195&214\\
 4:5&50&225&210\\
 2:3&50&250&220\\
 1:2&50&268&280\\
 1:3&50&330&340\\
 1:10&50&610&620
\end{array}
\]
The nominal ratios summarize an independence calculation.
Let \(n_2,n_3\) be the displayed case counts for \(D_2,D_3\).
Using these frequencies exactly gives
\(R=37500/(n_2n_3)\), not an exact equality to every ratio
label. At the nominal \(2:3\) condition, for example,
\(R=15/22\). The stimuli describe frequencies as approximate.
Three symptom sets, disease names and orders were counterbalanced,
and answer order varied randomly. Page two requested a
justification and a separate joint-probability mathematics problem.
The paper does not print that mathematics problem's full wording
\cite[pp.~239--240, table~1]{Lombrozo2007SimplicityProbability}.

Note 3 reports an auxiliary study of twenty undergraduates asked
whether someone with one disease was more or less likely to have
another. Half chose each direction. The question offered no
equal-likelihood option. Equal aggregate counts in a forced
directional choice do not establish that individuals believe
the diseases independent; they cannot supply the independence
assumption used in the benchmark.

\subsection{Choices and the fitted interpretation}

Overall, 71 percent chose the simpler explanation and 29 percent
the focal complex alternative. The reported omnibus comparison was
\(\chi^2(7)=25,p<.01\).
Figure 1 shows declining simple-choice proportions as the
complex explanation becomes more prevalent, with local departures
from monotonicity. The source gives 94 percent simple choices
at \(15:1\), 89 at \(9:10\), 61 at \(2:3\), 67 at \(1:2\),
and 39 at \(1:10\) in its text and later comparisons.
The other plotted coordinates are not supplied as exact tabular
percentages. Figures 1 and 2 show error bars but their captions
and text do not define whether these are standard errors or
intervals; no interpretation is supplied by assumption.

The logistic analysis used the natural logarithm of the probability
ratio to predict simple-choice frequency. The printed coefficients
are slope .795, standard error .217, and constant \(-1.306\),
standard error .248. The source reports \(p<.01\) for the
constant and \(p=.8\) for comparison of the slope with one.
Nonrejection of a unit slope does not prove ideal probability use.
The source interprets the constant as an effective prior preference
of about 79 percent for simplicity, with a reported 95 percent
confidence interval of 69--86 percent, or an odds multiplier
about four, interval 2.3--6
\cite[pp.~241--242]{Lombrozo2007SimplicityProbability}.

\begin{hbdefinition}{def:p19-explanation-choice-curve}{An explicitly oriented response model}
For \(R>0\), real \(a\) and real \(b\), define a simple-answer
choice probability
\[
 q(R)=\frac{\exp(a)R^b}{1+\exp(a)R^b}.
\]
This is a model of response frequency, unless a separate generative
interpretation has been proved. Its log odds are \(a+b\log R\).
In this orientation, the source's reported 79 percent interpretation
corresponds to \(a=+1.306\), giving
\(\exp(a)\approx3.69\) and \(q(1)\approx.787\).
The negative constant printed in the article must not be inserted
unchanged into this positively oriented formula.
The paper does not supply its complete response coding and fitted
design matrix; this explicit orientation expresses its reported
interpretation rather than reconstructing an unreported fit.
\end{hbdefinition}

\begin{hbtheorem}{thm:p19-explanation-choice-model}{proved=annex:p19.explanation-choice-model}{Choice curves do not uniquely identify mental priors}
The displayed log-odds identity holds, and for \(b=1\),
\(q(R)>1/2\) exactly when \(\exp(a)R>1\).
In the independent fictional example with \(R=5/8\),
an editorial multiplier \(\exp(a)=4\) gives \(q=5/7\).
This choice-model value is neither of the event posteriors.

For every \(a\in\mathbb R,b>0\), the same aggregate curve can
be generated by a population of deterministic threshold choosers.
Consequently a fitted group curve does not uniquely establish
individual probability matching or one common mental prior.
Likewise, choosing the simple explanation and later reporting
a high prevalence can be associated when choice has no effect
on the report. A complete finite example appears in the annex.
\end{hbtheorem}

\subsection{Justifications, misunderstandings and mathematics}

Two coders classified all 144 justifications with 94 percent
agreement, resolving disagreements by discussion.
The new misunderstood category included interpreting a disease's
failure to cause a symptom as preventing that symptom.
Reported category percentages were 15 simple, 13 sufficient,
51 probable, 9 other and 12 misunderstood.
Simple, sufficient and misunderstood justifications occurred only
among simple choices and did not vary detectably across neighboring
baserate pairs (\(p>.2\)). Probability appeals supported both
choices and changed with baserate and choice (\(p<.01\)):
they increasingly supported the complex explanation.
Collapsing across choice left approximately half the people
in each paired-condition group appealing to probability,
15--20 out of 36.

The sensitivity analysis removed seventeen coded misunderstandings.
Its slope was .915, standard error .246, and its printed constant
\(-1.164\), standard error .261. The derived simple-preference
interpretation was 76 percent, reported 95 percent interval
66--84 percent. These parameters were reported not significantly
different from the main analysis; the source supplies no raw
records for independent refitting.

Sixty-nine percent answered the mathematics problem correctly,
with a reported comparison against chance of \(p<.01\).
The reported baserate comparison was
\(\chi^2(3)=.614,p=.893\); its full grouping table is not printed.
Correctness correlated with the two-cause choice at \(r=.094,p=.263\),
or \(r=.15,p=.122\) when only conditions favoring the complex
explanation were considered. These non-significant correlations
do not establish that mathematics, interpretation or joint-probability
beliefs played no role \cite[pp.~242--243, note~5]{Lombrozo2007SimplicityProbability}.

\section{Experiment 3: experienced cases and prevalence estimates}
\label{sec:p19-explanation-experiment-three}

The retained sample comprised 108 undergraduate and summer-school
students, 56 percent female, mean age 22, standard deviation 7,
participating for credit or a small gift.
Fourteen additional participants were excluded: two for multiple
answers, two for presentation of the wrong frequencies, and ten
for potential misunderstanding. A misunderstanding flag from
either one of the judges was sufficient for that exclusion.
This differs from retaining misunderstandings in Experiment 2's
main analysis.

Thirty-six PowerPoint presentations crossed four nominal prevalence
ratios, three disease orders arranged as a Latin square, and three
symptom sets. Three people saw each presentation, individually,
giving 27 per prevalence condition. Assignment was random.
The nominal conditions were \(15:1,9:10,1:2,1:10\).

Participants were told they coordinated an alien disease-research
team and should attend to prevalence. After learning the symptom
relations, which also remained on a response sheet, they saw a
diagnostic sign for the first disease and ten screens containing
75 aliens in total, four to eleven per screen. Screens lasted
two seconds with one-second gaps. This was repeated for each
of the other two diseases. The signs were red versus blue forehead
dots, purple smoke versus no smoke, and yellow versus black eyes
under a light. Sign order was fixed while disease-to-sign assignment
varied with disease order.
Because tests occurred in separate blocks, no disease-co-occurrence
record was supplied.

After these cases, participants chose the most satisfying explanation
from the same six-option structure, justified it, and finally
estimated each disease's population percentage.
The source says the frequencies corresponded to Experiment 2
with 75 rather than 750 cases, but does not tabulate the integer
positive counts in every presentation. Dividing every earlier
count by ten would create fractional cases, so no such constructed
counts are reported as observations
\cite[pp.~245--246]{Lombrozo2007SimplicityProbability}.

\subsection{Choice proportions and a reported-test discrepancy}

Simple-choice proportions were 96, 59, 56 and 41 percent at
the four nominal ratios, respectively, with 63 percent overall.
The printed omnibus result is \(\chi^2(3)=7.26,p<.01\).
The following calculation specifies what follows if the printed
statistic and reference degrees of freedom are the intended inputs.

\begin{hbdefinition}{def:p19-explanation-tail}{The stated chi-square reference calculation}
Let \(X\) have the chi-square reference law with three degrees
of freedom, namely \(\operatorname{gamma}(3/2,1/2)\) in the
shape/rate convention of \cref{def:p15-gamma-density}.
For \(x>0\), put \(Q_3(x)=\mathbb P(X\geq x)\).
For \(u\in\mathbb R\), define
\[
 \operatorname{erfc}(u)=\frac2{\sqrt\pi}\int_u^\infty e^{-v^2}\,dv.
\]
This is the conventional upper-tail calculation for the printed
inputs, not a proof that a particular contingency-table statistic
has that reference distribution under the study design.
\end{hbdefinition}

\begin{hbtheorem}{thm:p19-explanation-chi-tail}{proved=annex:p19.explanation-chi-tail}{The probability implied by the printed inputs}
For \(x>0\),
\[
 Q_3(x)=\operatorname{erfc}\!\sqrt{x/2}
       +\sqrt{2x/\pi}\,e^{-x/2}.
\]
Consequently \(Q_3(7.26)\approx.06406\), and in particular
\(Q_3(7.26)>.01\).
\end{hbtheorem}

Thus the printed statistic and degrees of freedom do not yield
the printed \(p<.01\) under that convention. The actual observed
statistic or analysis cannot be repaired by inventing participant
counts from rounded percentages. The choice proportions remain
reported data, while the omnibus significance claim is unresolved.

The source's cross-experiment comparisons report, for Experiment 3
versus Experiment 2: 96 versus 94 percent at \(15:1\),
\(\chi^2(1)=.09,p=.77\); 56 versus 67 at \(1:2\),
\(\chi^2(1)=.56,p=.46\); 41 versus 39 at \(1:10\),
\(\chi^2(1)=.02,p=.90\); and 59 versus 89 at \(9:10\),
\(\chi^2(1)=4.6,p<.05\).
These comparisons change participants, format and exclusion policy,
so they do not by themselves isolate a causal format effect.

The second logistic fit gave slope .62, standard error .157,
and printed constant \(-.77\), standard error .24.
Its simple-preference interpretation was 68 percent, reported
95 percent interval 57--78 percent.
The slope difference from Experiment 2 had \(p>.3\), while
the constant difference had \(p<.05\).
Note 6 reports no significant difference from Experiment 2's
76 percent interpretation after excluding misunderstandings.
Disease-order effects on choice had \(p>.1\), and its interaction
with baserate had \(p>.3\)
\cite[pp.~246--247]{Lombrozo2007SimplicityProbability}.

\subsection{Justifications and selective prevalence estimates}

Two coders reviewed the 108 retained justifications, initially
agreeing on 96 percent. Reported categories were 8 percent
simple, 31 sufficient, 56 probable and 5 other.
Neighboring ratios were combined for analysis.
The first, second and last categories almost always supported
the simpler choice. Among probability appeals, 66 percent
supported simplicity in the two more favorable ratio conditions,
while 81 percent supported complexity in the two less favorable
conditions. The overall frequency of probability appeals stayed
about 56 percent across base rates; no significant overall
justification-frequency effect was reported.

Figure 3 plots three separate prevalence-estimation panels,
one for each disease. Solid reference lines show the presented
prevalences; estimates are split by chosen explanation.
The first panel has a different vertical range from the other
two. The central numerical comparisons are:
\[
\begin{array}{c|rr|l}
 \text{Nominal ratio}&
 \text{Chose }D_1&\text{Chose }D_2,D_3&
 \text{Reported unequal-variance test}\\\hline
 1:2&19\%&8\%&t(14.6)=2.18,\ p<.05\\
 1:10&29\%&10\%&t(10.6)=1.74,\ p=.11
\end{array}
\]
These are estimates of \(D_1\)'s prevalence, not explanation-choice
percentages. Both tests were two-tailed. At \(15:1\) and \(9:10\)
the corresponding choice-group difference was not detected.
Variance was reported greater among the simple choosers in
the two unfavorable conditions.
The combined two-by-two analysis gave
\(F(1,50)=7.93,p<.01\) for explanation choice and
\(F(1,50)=.442,p=.51\) for its interaction with baserate.
The analogous \(D_2,D_3\) comparisons had \(p>.2\).
Order effects on prevalence estimates and their interactions
had \(p>.3\) \cite[pp.~248--249]{Lombrozo2007SimplicityProbability}.

The author argues against a general frequency-estimation deficit
and confusion of \(D_1\) with another disease: errors were selective,
and only one participant estimated \(D_1\) above either other disease
in the \(1:2\) or \(1:10\) conditions.
The article acknowledges both causal directions, then favors
explanatory commitment changing estimates. But explanation choice
was not randomized. A participant's pre-existing internal estimate
could influence both the choice and the later report.
The finite construction in the annex shows why eliciting the report
after the choice does not eliminate that possibility.
For example, give four fictional people pre-choice estimates
8, 8, 20 and 20 percent. Let each choose the simple explanation
exactly when that estimate exceeds a fixed 15 percent comparison
value, then report the same estimate afterward. The choice groups
now report means 20 and 8 percent, although making a different
choice would not change anyone's estimate in this construction.

\section{Experiment 4: explicit conjunction frequencies and a changed question}
\label{sec:p19-explanation-experiment-four}

Forty-eight undergraduate and summer-school students from an
unnamed elite university participated, 58 percent female, mean
age 20, standard deviation 3, for credit, payment or a small gift.
One additional non-native English speaker was excluded for failing
to understand the task.
Participants were randomly assigned to two conditions; separate
condition counts are not printed in the method.
Both used the nominal \(2:3\) scenario.

The joint-probability condition retained the satisfying-explanation
question and added pairwise disease frequencies. Its example gave
about seventeen Brom--Stemmel cases, fifteen Brom--Pilt cases,
and seventy-three Stemmel--Pilt cases. Thus the focal complex
explanation's joint count was supplied directly rather than
requiring multiplication of marginals.
The likely condition retained only the individual disease frequencies
but asked for the \emph{most likely} explanation.
The second questionnaire page, including justification and the
mathematics question, matched Experiment 2.
Disease order, symptoms and answer order were varied as described
in that source \cite[pp.~250--251]{Lombrozo2007SimplicityProbability}.

Only 17 percent chose the simple explanation with joint frequencies,
compared with 46 percent under the likely wording.
Their reported difference was \(\chi^2(1)=4.75,p<.05\).
Against Experiment 2's 61 percent at the same nominal ratio,
the joint condition gave \(\chi^2(1)=8.85,p<.01\), and the
likely condition \(\chi^2(1)=1.34,p>.2\).
The within-experiment contrast therefore differs from merely
asking a likelihood question. It does not exclude every possible
task-demand account; the two manipulations also change available
information, and this is not a complete factorial isolation of
all wording and information effects.

Two justification coders agreed initially on 88 percent and
resolved disagreements by discussion.
Simple choices were justified 13 percent by simplicity,
40 by sufficiency, 33 by probability and 13 by other reasons.
Complex choices were justified 94 percent by probability and
6 by other reasons. No significant condition difference in
these justifications was reported.
Seventy-five percent answered the math problem correctly;
the condition comparison was \(t(46)=-.656,p=.525\).
Its correlation with explanation choice was \(r=.05,p=.83\)
in the joint condition and \(r=.51,p<.05\) in the likely condition.
In the latter, simple choosers were more likely to answer the
math question incorrectly
\cite[pp.~251--252]{Lombrozo2007SimplicityProbability}.

\section{A complete fictional calculation and its interpretation}
\label{sec:p19-explanation-worked}

Return to the explicitly independent fictional prevalences
\((.1,.4,.4)\). The observation event has probability
\(.1+.16-.016=.244\).
The one-cause event has posterior \(25/61\), whereas the
two-cause event has posterior \(40/61\).
They overlap, so these two posteriors sum to \(65/61\), not one.
Their ratio is \(5/8\).
The normalized two-label score is \(5/13\);
adding the stipulated choice multiplier four yields \(5/7\).
These four different numerical summaries have four explicitly
different meanings. None is a new participant observation.

The logical minimum-cardinality criterion selects \(\{D_1\}\)
in this model, while ranking inclusive explanatory events by
posterior probability favors \(D_2\cap D_3\).
Changing their unreported dependence can reverse that comparison
without changing any marginal prevalence.
This is a concrete failure boundary for treating cause count
as a universal probability rule.

\begin{hbmethod}{def:p19-explanation-assessment}{Assessing an explanatory-choice claim}
First specify whether the question concerns sufficiency, satisfaction,
probability or a decision with a stated objective.
List the candidate explanations and their inclusive or exclusive
meanings. State which observations they cover and which background
premises establish that coverage.
For a probability comparison, identify the population law, the
observation or selection event, the marginal and joint information,
and each independence assumption.
Compute event posteriors separately from any normalized answer
weights or response-frequency model.
Finally report the observed choice measure, analysis sample,
exclusions and uncertainty before interpreting an internal process.
An analogy to another explanatory task supplies a proposed
correspondence to examine, not an automatic transfer of its evidence.
This procedure is a handbook synthesis of the distinctions above.
\end{hbmethod}

\section{Evidence boundaries and links to other reasoning tasks}
\label{sec:p19-explanation-evidence-boundary}

The four main experiments support a role for both cause count and
provided probability information in these participants' selections.
They do not identify one universal subjective prior, prove that
the world is simple, or establish a unique Bayesian cognitive
implementation. The source's group logistic interpretation is one
account of the choice data; the mathematics above states the
additional event, response and dependence assumptions it requires.

The source distinguishes two broad accounts: explanatory quality
may affect perceived probability, or simplicity may be favored
because the world is believed simple without mediation by
explanatory quality. The experiments do not separate those accounts
fully. They also concern causal \emph{tokens}: which known causes
are present in one alien. Parsimony over kinds of entities or laws,
or over parameters in curve fitting, is a different problem.
No identity with minimum description length, AIC, Kolmogorov
complexity or Bayesian model integration is established here.

Recruitment institution and dates, complete participant records,
the math-question text, full analysis code and the definitions
of the graphical error bars are not supplied.
Selected results were presented at conferences in 2004; those
dates are not silently assigned to every experiment's recruitment.
Percentages and many uncertainty summaries are rounded.
Participant replacement, different misunderstanding rules and
post-choice subgroup comparisons limit the corresponding inferences.

The analogy experiments in
\cref{sec:p19-analogy-and-problem-transfer} investigate remembering,
noticing and adapting a supplied source when generating a solution.
Lombrozo's experiments instead supply the candidate explanations
and ask participants to select among them.
The distinction between a structurally available candidate and
logical truth in \cref{thm:p19-analogy-candidate-truth} is relevant;
it is not evidence that these different tasks share one measured
mechanism or interchangeable success rates.

A replication must state whether it preserves the one-versus-two
cause comparison, deterministic symptom descriptions, unfamiliar
material, offered menu, prevalence format, instruction, exclusion
policy and recruitment frame. Testing interacting causes, multiple
causes on both sides, explanation generation, new populations or
practical decisions would extend the evidence rather than merely
repeat the reported task. No universal explanatory prescription
follows from the chosen cases alone.
